Fix a realized graph \(A\), the preassignment response functions \((f_i)_{i\in\mathcal V_n}\), and \(h\in\mathcal H_{n,d}\).  All probabilities, expectations, and variances until stated otherwise are over the assignment.  By \({\rm ass:bernoulli\text{-}design}\), its conditional law is the product of \(n\) Bernoulli-one-half laws.  Put
\[
 X=Z-\tfrac12\mathbf1_n,\qquad
 P_A=\widehat\Pi_r(A),\qquad C_A=P_AA,
 \qquad F_h(Z)=\widehat\theta_h.
\tag{1}
\]
The definition of \(A\) gives \(A_{ii}=0\), and \({\rm def:empirical\text{-}projector}\) makes \(P_A\) an \(n\)-dimensional orthogonal spectral projector even when \(r>n\).

We first establish the coordinate oscillations.  Let \(z^{(j)}\) differ from \(z\) only by flipping coordinate \(j\).  The finite maximum in \({\rm def:pc\text{-}handle}\) gives
\[
 |q_{i,h}(z)-q_{i,h}(z^{(j)})|\leq D_{ij,h}.
\tag{2}
\]
Indeed, if \(A_{ij}=0\), the treated-neighbor count \(M_i\) is unchanged and both sides of (2) vanish.  If \(A_{ij}=1\), then \(M_i\) moves between two consecutive integers in the maximum defining \(D_{ij,h}\).  Division by \(\mu_{2,i}(h)\) is legitimate for \(N_i>0\) by \({\rm lem:exact\text{-}score\text{-}moments}\), while the zero-degree conventions make (2) immediate for \(N_i=0\).  Since \(|X_j|=1/2\), the definitions of \(Q_{i,h}\), \(\Gamma_{i,h}\), and \(\widetilde\gamma_{i,h}\) imply
\[
 |\widetilde\gamma_{i,h}(z)|
 \leq Q_{i,h}+2\sum_{k\in\mathcal V_n}|(C_A)_{ik}|
 =\Gamma_{i,h}.
\tag{3}
\]
A flip changes \(X_j\) by one in absolute value, so (2) also gives
\[
 |\widetilde\gamma_{i,h}(z)-\widetilde\gamma_{i,h}(z^{(j)})|
 \leq D_{ij,h}+4|(C_A)_{ij}|.
\tag{4}
\]

The level part of \({\rm ass:response\text{-}envelope}\) yields \(|Y_i^{\mathrm{obs}}(z)|\leq B\).  If \(j=i\), flipping \(j\) changes only the own-treatment argument, and hence changes the response by at most \(2B\).  If \(j\ne i\) and \(A_{ij}=1\), then \(N_i>0\), and the exposure changes by \(1/N_i\).  The first-derivative part of \({\rm ass:response\text{-}envelope}\) and the fundamental theorem of calculus give
\[
 |f_i(z_i,x+N_i^{-1})-f_i(z_i,x)|
 =\left|\int_x^{x+N_i^{-1}}\partial_t f_i(z_i,t)\,dt\right|
 \leq B/N_i.
\tag{5}
\]
In all other cases \({\rm ass:anonymous\text{-}interference}\) makes the response unchanged.  With the convention in \({\rm def:pc\text{-}handle}\),
\[
 |Y_i^{\mathrm{obs}}(z)-Y_i^{\mathrm{obs}}(z^{(j)})|
 \leq B\left\{2\mathbf1\{i=j\}+\frac{A_{ij}}{N_i}\right\}.
\tag{6}
\]
Applying \(|gy-g'y'|\leq|g-g'||y|+|g'||y-y'|\) to (3), (4), and (6) proves
\[
 |\widetilde\gamma_{i,h}(z)Y_i^{\mathrm{obs}}(z)
 -\widetilde\gamma_{i,h}(z^{(j)})Y_i^{\mathrm{obs}}(z^{(j)})|
 \leq\beta_{ij,h}.
\tag{7}
\]
Summing (7) and using \({\rm def:estimator}\) therefore gives
\[
 |F_h(z)-F_h(z^{(j)})|
 \leq c_{j,h}:=\frac1n\sum_{i\in\mathcal V_n}\beta_{ij,h}.
\tag{8}
\]
By the definition of \(\widehat V_{\mathrm{pc},h}\) and \({\rm def:radii}\),
\[
 \sum_{j\in\mathcal V_n}c_{j,h}^2
 =\frac1{n^2}\sum_j\left(\sum_i\beta_{ij,h}\right)^2
 =4\widehat V_{\mathrm{pc},h}=4\overline s_h(A)^2.
\tag{9}
\]

For completeness, the same oscillations give the stated variance certificate with its sharper factor.  Let \(\mathcal F_j=\sigma(Z_1,\ldots,Z_j)\), \(M_j=\mathbb E_Z(F_h\mid\mathcal F_j)\), and \(E_j=M_j-M_{j-1}\).  The tower property gives \(\mathbb E_Z(E_jE_k)=0\) for \(j<k\), and therefore
\[
 \operatorname{Var}_Z(F_h)=\sum_{j=1}^n\mathbb E_ZE_j^2.
\tag{10}
\]
Conditional on \(\mathcal F_{j-1}\), the two possible values of \(M_j\) differ by an average of differences bounded by \(c_{j,h}\).  They have equal conditional mass by \({\rm ass:bernoulli\text{-}design}\), so their conditional variance is at most \(c_{j,h}^2/4\).  Equations (9)--(10) yield
\[
 \operatorname{Var}_Z(\widehat\theta_h\mid A,(f_i)_i)
 \leq\frac14\sum_jc_{j,h}^2
 =\widehat V_{\mathrm{pc},h}=\overline s_h(A)^2.
\tag{11}
\]

We next prove the bias certificate.  For \(N_i>0\), the zero diagonal implies that \(q_{i,h}\) depends only on neighbor assignments and hence is independent of \(Z_i\).  Conditioning first on \(Z_i=a\), using \({\rm ass:anonymous\text{-}interference}\) and \({\rm lem:exact\text{-}score\text{-}moments}\), and then averaging the two equiprobable values of \(a\) gives
\[
 \left|\mathbb E_Z[q_{i,h}Y_i^{\mathrm{obs}}]
 -\frac{\partial_xf_i(1,1/2)+\partial_xf_i(0,1/2)}2\right|
 \leq\frac L6h_i(h)^2.
\tag{12}
\]
For \(N_i=0\), the score is zero, while the first-derivative part of \({\rm ass:response\text{-}envelope}\) bounds the absolute target summand by \(B\).  The score component therefore has conditional bias at most
\[
 \frac{L}{6n}\sum_{i:N_i>0}h_i(h)^2
 +\frac Bn\sum_i\mathbf1\{N_i=0\}.
\tag{13}
\]

It remains to control the PC subtraction.  Centering and \(|Y_i^{\mathrm{obs}}|\leq B\) imply
\[
 |\mathbb E_Z[X_iY_i^{\mathrm{obs}}]|\leq B/2.
\tag{14}
\]
For \(j\ne i\) with \(A_{ij}=1\), condition on all coordinates except \(Z_j\).  Equation (5) and the Bernoulli-one-half calculation imply
\[
 |\mathbb E_Z[X_jY_i^{\mathrm{obs}}]|
 =\frac14\left|\mathbb E_Z[Y_i^{\mathrm{obs}}\mid Z_j=1]
 -\mathbb E_Z[Y_i^{\mathrm{obs}}\mid Z_j=0]\right|
 \leq\frac B{4N_i}.
\tag{15}
\]
For \(j\ne i\) with \(A_{ij}=0\), anonymous interference and assignment independence make this expectation zero.  Consequently,
\[
\begin{aligned}
 \frac4n\sum_i\left|\sum_j(C_A)_{ij}\mathbb E_Z[X_jY_i^{\mathrm{obs}}]\right|
 &\leq\frac{2B}{n}\sum_i|(C_A)_{ii}|\\
 &\quad+\frac Bn\sum_{i:N_i>0}\frac1{N_i}
       \sum_{j:A_{ij}=1}|(C_A)_{ij}|\\
 &=\widehat\Delta_{\mathrm{pc},h}.
\end{aligned}
\tag{16}
\]
Combining (13)--(16) with \({\rm def:target}\), \({\rm def:estimator}\), and \({\rm def:radii}\) proves
\[
 \left|\mathbb E_Z(\widehat\theta_h-\theta_n\mid A,(f_i)_i)\right|
 \leq\overline b_h(A).
\tag{17}
\]

We also verify the graph-norm simplification.  Since \(P_A\) is a spectral projector of the symmetric matrix \(A\), the matrices commute and \(C_A=P_AA\) is symmetric.  In a spectral decomposition \(C_A=\sum_t\nu_tu_tu_t^\top\), orthonormality yields
\[
 \sum_i|(C_A)_{ii}|
 \leq\sum_t|\nu_t|\sum_i u_{t,i}^2
 =\sum_t|\nu_t|=\|C_A\|_*.
\tag{18}
\]
Cauchy--Schwarz within every positive-degree row and then across at most \(n\) rows gives
\[
\begin{aligned}
 \sum_{i:N_i>0}\frac1{N_i}\sum_{j:A_{ij}=1}|(C_A)_{ij}|
 &\leq\sum_{i:N_i>0}
 \left\{\frac1{N_i}\sum_{j:A_{ij}=1}(C_A)_{ij}^2\right\}^{1/2}\\
 &\leq\sqrt{nG_A(C_A)}.
\end{aligned}
\tag{19}
\]
Substitution of (18)--(19) in \({\rm def:pc\text{-}handle}\) proves
\[
 \widehat\Delta_{\mathrm{pc},h}
 \leq\frac{2B}{n}\|C_A\|_*+B\sqrt{\frac{G_A(C_A)}n}.
\tag{20}
\]

Apply \({\rm lem:mcdiarmid\text{-}bounded\text{-}differences}\) to the conditional product law and the oscillations (8).  If \(\overline s_h(A)>0\), set
\[
 t_h=\overline s_h(A)\sqrt{2\log(2/\alpha)}.
\]
Using (9) and \(0<\alpha<1/2\),
\[
\begin{aligned}
 \Pr_Z\{|F_h-\mathbb E_ZF_h|>t_h\mid A,(f_i)_i\}
 &\leq2\exp\left\{-\frac{2t_h^2}{\sum_jc_{j,h}^2}\right\}\\
 &=2\exp\{-\log(2/\alpha)\}=\alpha.
\end{aligned}
\tag{21}
\]
If \(\overline s_h(A)=0\), then (9) gives \(c_{j,h}=0\) for all \(j\).  Equation (8) then shows that \(F_h\) is unchanged by every single-coordinate flip.  The hypercube \(\{0,1\}^n\) is connected by such flips, so \(F_h\) is constant on the assignment support and hence \(F_h=\mathbb E_ZF_h\) almost surely; thus (21) remains valid.  Combining (17) and (21) by the triangle inequality yields
\[
 \Pr_Z\!\left\{
 |\widehat\theta_h-\theta_n|>
 \overline b_h(A)+\overline s_h(A)\sqrt{2\log(2/\alpha)}
 \mathrel{\Big|}A,(f_i)_i\right\}\leq\alpha.
\tag{22}
\]
The derivative envelope and \({\rm def:target}\) give \(|\theta_n|\leq B\), so intersecting with \([-B,B]\) cannot remove the target on the event in (22).  This proves the fixed-bandwidth conditional coverage in the corrected \({\rm def:honest\text{-}interval}\), and integration proves uniform unconditional coverage over the maintained model class.

We finally verify the selector and computational claims.  On \(\mathcal H_{n,d}\ne\varnothing\), \({\rm def:bandwidth\text{-}set}\) gives the compact interval \([\ell,u]\) with \(\ell=d^{-1}\) and \(u=\min\{d^{-1/2},1/2\}\).  The set \(\mathcal K(A)\) in \({\rm def:operational\text{-}bandwidth}\) contains both endpoints and at most
\[
 2+\sum_{i:N_i>0}(N_i+1)\leq2+n^2
\tag{23}
\]
candidates.  At every candidate, \({\rm lem:exact\text{-}score\text{-}moments}\) makes the positive-degree normalizers positive, so the finite maxima and the two radii are finite.  A minimum over this finite nonempty set therefore exists, and the smallest-minimizer rule makes \(\widehat h(A)\) single valued.  The graph sample space is the finite set of symmetric zero-diagonal binary matrices; hence every such single-valued graph map is Borel measurable.  Exact binomial sums compute \(\mu_{2,i}(h)\), finite enumeration computes \(Q_{i,h}\) and \(D_{ij,h}\), one eigendecomposition computes \(C_A\), and finite matrix arithmetic computes the remaining quantities.  Equation (23) therefore proves polynomial-time computation in the real-arithmetic model already used by the frozen construction.  Since \(\widehat h(A)\) depends on neither \(Z\) nor the outcomes, it is constant conditional on \(A,(f_i)_i\).  Applying (8)--(22) at that fixed conditional value proves every selected certificate and selected-coverage assertion without a union bound.

For a degree \(N_i\), the truncated score changes only when \(h\) crosses one of the finitely many distances \(|m/N_i-1/2|\).  Thus (23) also represents all distinct score regimes.  Finally, \({\rm lem:exact\text{-}score\text{-}hoeffding\text{-}projections}\) gives
\[
 q_h^{(1)}=4AX,\qquad q_h^{(2)}=0.
\]
Therefore the first projection of \(q_h-4P_AAX\) is \(4(I_n-P_A)AX\), while its second projection is zero; neither expression depends on \(h\).  Equations (8), (9), (11), (17), (20)--(23), and this projection calculation prove all parts of the proposed theorem without a sharp stochastic-rate assertion.